\begin{textsample}{NEG Dim 1 – 2011\_c4 – Score -4.00 – 2011\_c4/enade\_2011\_tec\_red\_comp\_c877d48bcc.txt}  \label{ex:f1_neg_115}
CAPÍTULO 4 PERCEPÇÃO DA PROVA As análises feitas neste capítulo tratam das percepções dos concluintes da Área de Tecnologia em Redes de Computadores sobre a prova aplicada no ENADE/2011.
Estas percepções foram mensuradas por meio de nove questões que avaliaram desde o grau de dificuldade da prova até o tempo gasto para concluí -la.
As percepções sobre a prova foram relacionadas com o desempenho dos estudantes e com a Grande Região de funcionamento do curso.
O questionário de percepção da prova encontra -se ao final do Anexo IV, que traz a reprodução da prova.
O desempenho dos estudantes foi classificado em quatro quartos.
Para tanto, esse desempenho foi ordenado de forma ascendente.
O percentil 25, P25, também conhecido como primeiro quartil, é a nota de desempenho que deixa um quarto ( 25 \% ) dos valores observados abaixo e três quartos acima.
A Figura 1 apresenta uma ilustração deste conceito.
O quarto inferior de desempenho é composto pelas notas abaixo do primeiro quartil.
Já o percentil 75, P75, também conhecido como terceiro quartil, é o valor para o qual há três quartos ( 75 \% ) dos dados abaixo e um quarto acima dele.
O quarto superior de desempenho é composto pelas notas iguais ou acima do terceiro quartil.
O percentil 50, P50, também conhecido como mediana, é o valor que divide as notas em dois conjuntos de igual tamanho.
O segundo quarto inclui valores entre o primeiro quartil ( P25 ) e a mediana.
O terceiro quarto contém os valores entre a mediana ( P50 ) e o terceiro quartil ( P75 ).
Vale ressaltar que percentis, quartis e medianas são pontos que não obrigatoriamente pertencem ao conjunto original de dados, ao passo que os quartos são subconjuntos dos dados originais.
A seguir, serão apresentados gráficos com resultados selecionados, relativos às nove questões avaliadas por grupos de estudantes.
Os gráficos apresentam nas barras o percentual de alunos que assinalaram uma das opções ou a soma das porcentagens daqueles que assinalaram duas ( ou três ) delas.
Por exemplo, para as questões 1 e 2, os gráficos apresentam a porcentagem total de participantes que assinalaram as opções ( D ) difícil e ( E ) muito difícil.
Em \textbf{cada} barra foram assinalados também os extremos do intervalo de confiança de 95 \% como linhas verticais unidas por uma linha horizontal na forma da letra H maiúscula, semelhantemente aos gráficos do Capítulo 3.
As Tabelas no Anexo II apresentam os valores absolutos e a distribuição percentual das alternativas válidas das nove questões, segundo o mesmo recorte de desempenho dos alunos e Grande Região de funcionamento do curso. 4.1 GRAU DE DIFICULDADE DA PROVA 4.1.1 Componente de Formação Geral Ao avaliarem “ Qual o grau de dificuldade desta prova na parte de Formação Geral? ” ( Questão 1 ), 28,4 \% do grupo de inscritos e presentes optaram pelas alternativas difícil ou muito difícil.
Entretanto, para mais da metade dos estudantes ( 58,6 \% ), o Componente de Formação Geral da prova foi considerado com grau de dificuldade médio ( Gráfico 4.1, Gráfico 4.2 e, no Anexo II, a Tabela II.1 ).
O percentual de estudantes que consideraram a prova como difícil ou muito difícil foi maior na região Sudeste, onde a proporção foi de 30,3 \%, enquanto a de menor incidência foi a região Nordeste, 21,6 \%.
No Gráfico 4.1 que não há diferenças estatisticamente significativas entre estas duas regiões.
Nas Grandes Regiões, a proporção de presentes à prova que consideraram o Componente de Formação Geral como sendo de grau de dificuldade médio \textbf{esteve} entre 56,3 \% na região Sudeste e 61,7 \% na região Nordeste.
O percentual de alunos que consideraram a prova difícil ou muito difícil foi mais elevado no quarto de pior desempenho ( 33,0 \% ), e o menor valor é observado no quarto de melhor desempenho ( 23,6 \% ).
Para todos os quartos de desempenho a alternativa modal para esta pergunta foi médio, sendo 56,4 \% no 1º quarto a 57,5 \% no 4º quarto. 4.1.2 Componente de Conhecimento Específico Ao responderem à Questão 2 — “ Qual o grau de dificuldade desta prova na parte de Componente Específico? ” — 53,0 \% do grupo de estudantes classificaram -na como difícil ou muito difícil.
Além disso, o Componente de Conhecimento Específico da prova foi considerado com grau de dificuldade médio por 41,9 \% dos alunos ( Gráfico 4.3, Gráfico 4.4, e, no Anexo II, a Tabela II.2 ).
A análise das respostas dos estudantes quanto ao grau de dificuldade do Componente de Conhecimento Específico da prova, agregado por Grande Região, mostra que a diferença entre a maior e a menor proporção de alunos que a avaliaram como difícil ou muito difícil é estatisticamente significativa : 57,4 \% na região Centro-Oeste e 44,6 \% na Nordeste.
O percentual de alunos que classificaram o grau de dificuldade como médio, no Componente de Conhecimento Específico, variou de 37,9 \% a 49,7 \%, para regiões Centro-Oeste e Nordeste, respectivamente.
Considerando -se a avaliação da dificuldade das questões do Componente de Conhecimento Específico da prova, de acordo com o desempenho dos estudantes, não se observa diferença estatisticamente significativa.
Em todos os quartos, a proporção dos que classificaram a parte específica como difícil ou muito difícil variou de 50,2 \% ( 1º quarto ) a 56,6 \% ( 3º quarto ).
No outro extremo das opções, a proporção dos estudantes que classificaram a parte específica como fácil ou muito fácil foi bastante reduzida, variou de 3,8 \% ( 3º quarto ) a 6,6 \% ( 1º quarto ). 4.2 EXTENSÃO DA PROVA EM RELAÇÃO AO TEMPO TOTAL Indagados quanto à extensão da prova, em relação ao tempo total oferecido para a sua resolução ( Questão 3 ), os estudantes apontaram, com maior incidência, a alternativa que considerava a extensão adequada, para todas as agregações consideradas ( Gráfico 4.5, Gráfico 4.6, e, no Anexo II, a Tabela II.3 ).
O percentual de alunos que responderam ser a extensão da prova adequada foi de 61,2 \%.
Já 33,7 \% dos inscritos e presentes consideraram que o exame foi longo ou muito longo e um pouco mais de 5 \% o avaliaram como curto ou muito curto.
Entre as Grandes Regiões a proporção daqueles que avaliaram a prova como longa ou muito longa em relação ao tempo total destinado à sua resolução variou pouco : de 30,0 \% na região Norte a 35,9 \% na região Sudeste.
As diferenças entre as regiões não são estatisticamente significativa.
Considerando -se o desempenho dos alunos, observa -se que no 1º quarto 58,0 \% dos estudantes consideraram a extensão da prova adequada e, no 3º quarto, a proporção dos que emitiram tal opinião foi a mais elevada, 64,5 \%.
No 2º quarto e no 4º quarto a proporção de estudantes que consideraram a extensão da prova adequada foi praticamente a mesma, 61,2 \% e 61,1 \%, respectivamente.
No Gráfico 4.6, pode -se constatar que para os quatro quartos de desempenho a proporção de estudantes que consideraram a prova longa ou muito longa, em relação ao tempo total destinado à sua resolução, manteve -se em \textbf{torno} de 33,7 \%.
Não há diferenças estatisticamente significativas nem tendência de crescimento ou de decrescimento associada ao aumento de desempenho. 4.3 COMPREENSÃO DOS ENUNCIADOS DAS QUESTÕES 4.3.1 Componente de Formação Geral Com relação aos enunciados das questões do Componente de Formação Geral ( Questão 4 ), as opiniões foram positivas, já que 70,1 \% dos alunos avaliados consideraram os enunciados de todas ou da maioria das questões claros e objetivos ( Gráfico 4.7, Gráfico 4.8, e, no Anexo II, a Tabela II.4 ).
Na análise regional, a percentagem de estudantes que avaliaram que todos ou a maioria dos enunciados das questões do Componente de Formação Geral \textbf{estavam} claros e objetivos variou de 65,9 \%, na região Sudeste, a 77,8 \% na região Nordeste, com diferença estatisticamente significativa entre essas duas regiões.
A análise das percepções dos estudantes sobre a clareza e objetividade dos enunciados permite afirmar que todos ou a maioria dos enunciados de questões relativas ao Componente de Formação Geral foram considerados claros e objetivos para a maior parte dos respondentes ( maior do que 65 \% em todas as regiões e maior do que 62 \% para todos os quartos de desempenho ).
Na análise por desempenho observa -se que a proporção dos que consideraram todos ou da maioria dos enunciados das questões claros e objetivos cresce conforme o desempenho aumenta, com diferenças estatisticamente significativas entre o 1º e o 4º quarto.
No quarto superior, a clareza e objetividade de todos ou da maioria dos enunciados das questões foi percebida por 74,9 \% dos alunos e no inferior tal avaliação foi emitida por 62,7 \% deles. 4.3.2 Componente de Conhecimento Específico Com relação aos enunciados das questões do Componente de Conhecimento Específico da prova, para 66,2 \% dos estudantes avaliados da Área de Tecnologia em Redes de Computadores a clareza e a objetividade ( Questão 5 ) \textbf{estavam} presentes em todas ou na maioria das questões ( Gráfico 4.9, Gráfico 4.10, e no Anexo II, a Tabela II.5 ).
Em todas as Grandes Regiões brasileiras os enunciados foram considerados claros e objetivos para todas ou para a maioria das questões do Componente de Conhecimento Específico por mais do que 63 \% dos estudantes.
A diferença entre as proporções das regiões Sudeste e Nordeste, respectivamente, 63,0 \% e 72,3 \%, é estatisticamente significativa.
As proporções de estudantes que consideraram os enunciados das questões claros e objetivos assumem praticamente o mesmo valor nos quartos centrais : 67,9 \% no 2º quarto e 67,5 \% no 3º quarto.
Os quartos extremos apresentam a menor proporção, 60,6 \% no 1º quarto, e a maior proporção, 68,8 \% no 4º quarto. 4.4 SUFICIÊNCIA DAS INFORMAÇÕES/INSTRUÇÕES FORNECIDAS Ao avaliarem as informações/instruções fornecidas para a resolução das questões ( Questão 6 ), 79,9 \% dos respondentes da Área de Tecnologia em Redes de Computadores de todo o Brasil afirmaram que estas eram até excessivas ou suficientes em todas ou na maioria das questões ( Gráfico 4.11, Gráfico 4.12, e, no Anexo II, a Tabela II.6 ).
Quanto à distribuição de respondentes pelas Grandes Regiões observa -se que a proporção de estudantes que consideraram as informações/instruções fornecidas até excessivas ou suficientes em todas ou na maioria das questões foi sempre superior a 77 \%, chegando a quase 83,6 \% na região Sul.
As diferenças entre as regiões não chegaram a ser estatisticamente significativas.
Como mostra o Gráfico 4.12, o percentual de participantes que avaliaram as informações/instruções como até excessivas ou suficientes em todas ou na maioria das questões foi mais elevado no quarto superior ( 85,5 \% ), percentual superior à média nacional ( 79,9 \% ).
Já no quarto inferior, a suficiência das informações/instruções declarada como até excessiva em todas ou na maioria das questões foi percebida por 75,0 \% dos respondentes.
As diferenças entre os quartos de desempenho não são estatisticamente significativas, com exceção do 4º quarto em relação aos demais. 4.5 DIFICULDADE ENCONTRADA AO RESPONDER À PROVA \textbf{Perguntados} sobre as dificuldades com as quais se depararam ao responder à prova ( Questão 7 ), 18,2 \% dos estudantes apontaram o desconhecimento do conteúdo.
Para 49,7 \%, a forma diferente de abordagem do conteúdo foi indicada como dificuldade.
Já a falta de motivação para fazer a prova foi a dificuldade apontada por 17,3 \% dos respondentes.
Considerando -se todo o Brasil, 12,0 \% dos respondentes afirmaram que não tiveram qualquer tipo de dificuldade para responder à prova ( Tabela II.7 no Anexo II ).
Os Gráficos 4.13 e 4.14 apresentam os percentuais de estudantes que apontaram o desconhecimento do conteúdo como dificuldade percebida ao responder à prova.
Na análise por Grandes Regiões, o percentual de inscritos e presentes que apontaram o desconhecimento do conteúdo como dificuldade ao responder à prova não superou 22 \%.
Os percentuais variaram de 12,7 \% na região Norte a 21,2 \% na Centro-Oeste, sendo esta diferença significativa estatisticamente.
A forma diferente de abordagem do conteúdo foi a escolha modal dos estudantes, com percentuais que variaram de 46,7 \% ( região Nordeste ) a 57,2 \% ( Norte ).
O percentual de alunos que citou a falta de motivação como dificuldade variou de 12,7 \% ( região Centro-Oeste ) a 21,8 \% ( região Nordeste ).
Os que declararam não ter qualquer dificuldade para responder à prova variaram de 10,7 \% na região Centro-Oeste a 13,9 \% na Nordeste.
Com relação aos quartos de desempenho, o desconhecimento do conteúdo, foi a opção escolhida por 18,6 \% dos estudantes do quarto superior e por 15,6 \% do quarto inferior.
A alternativa modal para os alunos, quando agregados pelos quartos de desempenho, foi a dificuldade forma diferente de abordagem do conteúdo : 47,1 \% no quarto inferior e 52,3 \% do quarto superior assim o responderam. 4.6 CONTEÚDOS DAS QUESTÕES OBJETIVAS DA PROVA Ao analisarem os conteúdos das questões objetivas da prova ( Questão 8 ), um percentual pequeno dos estudantes avaliados, apenas 8,5 \%, afirmou que não estudou ainda a maioria desses conteúdos ( Gráficos 4.15, Gráfico 4.16, e a Tabelas II.8 no Anexo II ).
A maioria ( 52,0 \% ) afirmou ter estudado e aprendido muitos ou todos os conteúdos.
Na análise por Grande Região, a proporção de respondentes que escolheram a opção não estudou ainda a maioria desses conteúdos, foi pequena.
Observa -se que nas regiões Sudeste e Sul ( 10,0 \% ), apesar de pequenas, as proporções foram maiores do que a média nacional ( 8,5 \% ).
A diferença entre os percentuais das regiões Norte e Sudeste é estatisticamente significativa.
Mas quaisquer outras duas regiões não apresentam diferenças estatisticamente significativas.
Em todas as regiões, a maioria dos presentes afirmou ter estudado e aprendido muitos ou todos os conteúdos, com proporções variando entre 46,0 \% na região Norte e 60,2 \% na Nordeste.
Considerando -se separadamente as opiniões de estudantes dos quatro quartos de desempenho, observa -se que, no quarto inferior, 12,3 \% ofereceram como resposta que não estudou ainda a maioria desses conteúdos, sendo menos de 5,0 \% os do quarto superior com a mesma resposta.
A diferença entre os alunos que optaram por este motivo de dificuldade nos quartos extremos é estatisticamente significativa.
Conforme aumenta o desempenho há uma tendência decrescente na proporção dos que justificaram assim a dificuldade encontrada.
Tendo em conta o quarto superior, 69,1 \% dos alunos afirmaram ter estudado e aprendido muitos ou todos os conteúdos. 4.7 TEMPO GASTO PARA CONCLUIR A PROVA Ao responderem sobre o tempo de conclusão da prova ( Questão 9 ), mais da metade dos estudantes ( 58,4 \% ) afirmou ter gasto entre duas e quatro horas ( Gráfico 4.17, Gráfico 4.18 e, no Anexo II, a Tabela II.9 ).
Considerando -se as cinco Grandes Regiões brasileiras, na região Norte a proporção ( 74,8 \% ) dos que utilizaram entre duas e quatro horas para finalizar a prova foi a mais acentuada.
Nas demais regiões o percentual de alunos que dispensaram entre duas e quatro horas para concluir a prova ficou entre 50 \% e 70 \%.
Como mostra o Gráfico 4.17, há diferenças estatisticamente significativas entre as proporções das regiões Sul, Sudeste e Nordeste em relação à região Norte.
Também se observa diferença significativa estatisticamente entre a região Centro-Oeste e as proporções das regiões Sul e Sudeste.
Uma vez analisadas as alternativas escolhidas pelos estudantes dos diferentes quartos de desempenho, observa -se que uma maior proporção de participantes do quarto superior declarou ter gasto entre duas e quatro horas para concluir a prova quando comparados com os do quarto inferior, respectivamente 72,0 \% e 43,4 \%, consequência de uma tendência crescente das proporções com os quartos de desempenho.

% matched lemmas: cada, estar, perguntar, torno
\end{textsample}
